\section{应用与推广}\label{sec:applications}

\subsection{素数分布}

黎曼猜想对素数分布的刻画最为直接。在黎曼猜想下，
\begin{equation}
    \pi(x) = \operatorname{Li}(x) + O(\sqrt{x} \log x), \qquad
    p_n = \operatorname{Li}^{-1}(n) + O(\sqrt{n} \log^3 n),
\end{equation}
即第 $n$ 个素数与 $n \log n$ 的偏差被控制在平方根量级。舍恩菲尔德（Schoenfeld）在1976年给出了显式常数\cite{schoenfeld1976sharper}：在黎曼猜想下，对 $x \ge 2657$ 有
\begin{equation}
    |\pi(x) - \operatorname{Li}(x)| < \frac{1}{8\pi} \sqrt{x} \log x.
\end{equation}

\subsection{计算复杂度与密码学}

\subsubsection{素性判定}

米勒（Miller, 1976）提出的素性判定算法在广义黎曼假设下是确定性多项式时间的\cite{miller1976riemann}：对被测整数 $n$，只需对 $O(\log^2 n)$ 个基做 Miller--Rabin 型测试即可。无条件的确定性多项式判定直到2002年才由 AKS 算法实现，但 GRH 版米勒判定在实践中仍然更快。类似地，在 GRH 下可以在多项式时间内找到模 $p$ 的原根、构造有限域等。

\subsubsection{密码学中的素数}

RSA 等公钥密码体制依赖大素数的随机生成。素数间隔的估计（在黎曼猜想下可显著加强）保证了在长度为 $k$ 比特的区间内素数足够稠密，从而随机搜索在多项式时间内成功。虽然现行密码系统的安全性论证并不直接依赖黎曼猜想，但零点分布的任何意外（例如临界带外零点的存在）都会改变对最小非剩余、伪素数分布等量的界，从而影响安全参数的选取。

\subsection{$L$ 函数与代数数论}

GRH 是现代计算代数数论的基础性假设。在 GRH 下，类群、单位群、西罗子群等不变量的计算可以在多项式时间内严格完成（Bach 的工作）；Cohen 的经典教科书等普遍将“GRH 下的有效计算”视为标准设定。此外，GRH 蕴含关于欧拉积、类数渐近、椭圆曲线的 Selmer 群上界等一系列结果，是朗兰兹纲领中“标准猜想”的最低要求。

\subsection{数学物理}

零点与量子系统的联系是数学物理的经典话题。若希尔伯特—韦尔算子存在，则可构造一个哈密顿量，其能级恰为零点纵坐标，即“黎曼气体”。贝里—基廷提出 $xp$ 型哈密顿量\cite{berry1999heisenberg}；Connes 用阿代尔观点构造了截断的迹类算子，其非零谱恰为临界线上的零点\cite{connes1999trace}，尽管尚不能排除临界带内的谱。此外，$\zeta$ 函数的配分函数、统计力学中的李—杨定理都与零点分布存在形式上的类比。

\subsection{随机矩阵与数论的交汇}

GUE 统计在零点间距、矩、欧拉积等对象上的普遍出现，为“大规模数值实验—猜想—证明”的数论研究范式提供了范例。凯茨—萨纳克的对称类型分类\cite{katz1999zeros}不仅统一了函数域上的结果，也为研究 $L$ 函数族的一阶矩、矩分布与级联统计提供了正确模型，这一框架目前已被推广到椭圆曲线、模形式和自同构表示的 $L$ 函数族。
